\documentclass[12pt,letterpaper]{article}

% ── Paquetes ──────────────────────────────────────────────────────────────
\usepackage[utf8]{inputenc}
\usepackage[T1]{fontenc}
\usepackage[spanish,es-tabla]{babel}
\usepackage[margin=2.54cm]{geometry}
\usepackage{times}
\usepackage{setspace}
\usepackage{graphicx}
\usepackage{tabularx}
\usepackage{booktabs}
\usepackage{caption}
\usepackage{hyperref}
\usepackage{csquotes}
\usepackage{xcolor}
\usepackage{xspace}
\usepackage{amsmath}
\usepackage{amssymb}
\usepackage{enumitem}
\usepackage[pages=all]{background}
\usepackage[style=apa, backend=biber]{biblatex}

% ── Colores ────────────────────────────────────────────────────────────────
\definecolor{verdeETITC}{RGB}{0, 145, 60}

% ── Siglas ─────────────────────────────────────────────────────────────────
\newcommand{\LMS}{\textit{Learning Management System} (LMS)\xspace}
\newcommand{\ML}{\textit{Machine Learning} (ML)\xspace}
\newcommand{\RF}{\textit{Random Forest} (RF)\xspace}
\newcommand{\XGB}{\textit{XGBoost}\xspace}
\newcommand{\ETITC}{Escuela Tecnológica Instituto Técnico Central (ETITC)\xspace}
\newcommand{\PDI}{Plan de Desarrollo Institucional (PDI)\xspace}
\newcommand{\TyT}{Técnica y Tecnológica (TyT)\xspace}
\newcommand{\SPADIES}{Sistema de Prevención y Análisis de la Deserción en Instituciones de Educación Superior (SPADIES)\xspace}
\newcommand{\TRL}{\textit{Technology Readiness Level} (TRL)\xspace}
\newcommand{\ETL}{Extracción, Transformación y Carga (ETL)\xspace}
\newcommand{\CSV}{Valores Separados por Comas (CSV)\xspace}
\newcommand{\EDM}{\textit{Educational Data Mining} (EDM)\xspace}

% ── Marca de agua ─────────────────────────────────────────────────────────
\backgroundsetup{
  scale=10,
  color=verdeETITC,
  opacity=0.1,
  angle=0,
  contents={ETITC}
}

\addbibresource{referencias.bib}
\hypersetup{colorlinks=true, linkcolor=black, urlcolor=blue, citecolor=black}
\onehalfspacing

\begin{document}

% ────────────────────────────── PORTADA ──────────────────────────────────
\begin{titlepage}
  \centering
  \vspace*{1cm}
  {\Large \textbf{Sistema de analítica predictiva basado en variables de
  entorno dinámico para la detección temprana del riesgo de deserción:
  Caso jornada nocturna ETITC Sede Centro en Bogotá.}}\\[2cm]

  \textbf{Estudiantes:}\\
  Jhon Alejandro Correal Martínez\\
  Rafael Andrés Guzmán Rodríguez\\[1.5cm]

  \textbf{Docente:}\\
  Doris Constanza Alvarado Mariño\\[3cm]

  \vfill
  \ETITC\\
  Facultad de Sistemas – Proyecto de Investigación II\\
  Bogotá, Colombia, Mayo de 2026
\end{titlepage}

\newpage
\tableofcontents
\newpage

% ────────────────────────────── RESUMEN ──────────────────────────────────
\section{Resumen}

La deserción en la formación \TyT en Colombia afecta la movilidad social de miles de jóvenes cada año. Este proyecto diseñó y validó un sistema de analítica predictiva basado en \ML para detectar de forma temprana el riesgo de abandono en la jornada nocturna de la \ETITC Sede Centro. Dado que la institución no puede compartir microdatos reales de sus estudiantes por restricciones de privacidad establecidas en la Ley 1581 de 2012 \parencite{ley1581}, se construyó un conjunto de datos sintéticos de 1\,000 registros generado mediante reglas de negocio calibradas con los informes históricos de Bienestar Universitario de los últimos diez semestres \parencite{etitc_bienestar}. La simulación reprodujo la tasa de deserción histórica de la institución (17,6\,\%). Se entrenaron y compararon dos modelos supervisados: \RF y \XGB. El modelo \RF alcanzó una exactitud del 95,0\,\%, con precisión del 100\,\% y un área bajo la curva ROC de 0,976. El modelo \XGB, ajustado para reducir los casos no detectados, mejoró la sensibilidad al 85,7\,\%, identificando 30 de los 35 estudiantes en riesgo del conjunto de prueba. Los resultados confirman que las variables de autoevaluación y coevaluación del Corte~3, junto con la nota definitiva del mismo corte, son los predictores de mayor peso, lo que valida la arquitectura propuesta para una futura integración con datos reales anonimizados.

\vspace{0.3cm}
\noindent\textbf{Palabras clave:} deserción estudiantil; aprendizaje automático; datos sintéticos; minería de datos educativos; jornada nocturna.

\vspace{0.3cm}
\noindent\textbf{Ecuación de búsqueda:} \texttt{(deserción estudiantil AND machine learning AND educación técnica AND Colombia)}

\newpage

% ──────────────────────── FORMULACIÓN DEL PROBLEMA ───────────────────────
\section{Formulación del problema}

La deserción en la formación \TyT en Colombia es un problema que limita las oportunidades de vida de miles de personas. Según el \SPADIES 3.0, la tasa de deserción anual a nivel nacional se ubica en el 8,5\,\% \parencite{men2025}. La \ETITC registra una tasa del 12,05\,\%, que supera en 3,55 puntos porcentuales ese promedio \parencite{etitc2025b}. Para reducir esa brecha, el \PDI 2025-2032 de la institución fijó como meta bajar el indicador en 2,5 puntos porcentuales \parencite{etitc2025b}, apoyándose en el sistema Adviser V.11 de la empresa Bersoft y la plataforma RUSIA \parencite{bersoft2024}. Sin embargo, el análisis de los informes de caracterización de Bienestar Universitario de los últimos diez semestres muestra que estas herramientas no alcanzan a atender la situación real del estudiante nocturno \parencite{etitc_bienestar}.

El riesgo de abandono cambia semana a semana según lo que le ocurre al estudiante fuera del aula: la carga del trabajo, los problemas económicos, la fatiga acumulada. Los registros históricos de la institución revelan que entre el 67\,\% y el 70\,\% de los estudiantes nocturnos trabaja de manera activa, y que más del 52\,\% declara falta severa de tiempo para estudiar por cuenta propia \parencite{etitc_bienestar}. Esta combinación genera lo que algunos autores denominan ``pobreza de tiempo'' \parencite{arias2018}: una restricción crónica que afecta la regularidad académica antes de que ningún sistema de alertas la detecte.

A esto se suma una brecha técnica concreta. El sistema Adviser V.11 captura las notas oficiales, pero solo al cierre de cada corte \parencite{bersoft2024}. La plataforma RUSIA no integra los registros de entrega digital de Microsoft Teams. Las encuestas de Bienestar recogen datos socioeconómicos, pero estos no están conectados con las calificaciones \parencite{etitc_bienestar}. Cuando el sistema emite una alerta, el estudiante ya lleva semanas en riesgo sin que nadie lo haya notado.

Una limitación adicional es que la institución no puede compartir los microdatos individuales de los estudiantes para entrenar modelos de inteligencia artificial, debido a las restricciones de la Ley 1581 de 2012 \parencite{ley1581}. Esto obliga a desarrollar y validar el sistema con datos sintéticos antes de conectarlo con información real anonimizada.

\subsection{Pregunta problema}
¿De qué manera la integración de variables de caracterización socioeconómica con el rendimiento académico por cortes permite construir una arquitectura predictiva capaz de identificar de forma temprana el riesgo de deserción en la jornada nocturna de la ETITC Sede Centro?

\newpage

% ────────────────────────────── OBJETIVOS ────────────────────────────────
\section{Objetivos}

\subsection*{Objetivo general}

Prototipar y validar un modelo de analítica predictiva basado en \ML que integre datos de caracterización socioeconómica con el rendimiento académico por cortes, para la detección temprana del riesgo de deserción en la jornada nocturna de la ETITC Sede Centro en Bogotá.

\subsection*{Objetivos específicos}

\begin{enumerate}
  \item Identificar los factores socioeconómicos y académicos de mayor peso en la deserción estudiantil nocturna, a partir del análisis comparado de los informes históricos de Bienestar Universitario de los últimos diez semestres \parencite{etitc_bienestar}, con el fin de establecer las variables de entrada del modelo.
  \item Diseñar una arquitectura de datos relacional de tres fuentes independientes —caracterización socioeconómica, notas por corte y variable objetivo— que reproduzca la estructura fragmentada de los sistemas de información institucionales actuales, permitiendo integrar datos dispersos en un único conjunto de entrenamiento.
  \item Evaluar la consistencia predictiva de la arquitectura mediante el entrenamiento de dos algoritmos supervisados (\RF y \XGB) sobre un conjunto de datos sintéticos de 1\,000 registros, generado con reglas de negocio calibradas a las distribuciones históricas de la institución, midiendo exactitud, sensibilidad y AUC-ROC como indicadores de desempeño.
\end{enumerate}

\newpage

% ──────────────────────────── JUSTIFICACIÓN ──────────────────────────────
\section{Justificación}

Cada semestre, una parte significativa de los estudiantes nocturnos de ingeniería de la ETITC abandona sus estudios sin que la institución haya podido actuar a tiempo. Los informes de Bienestar Universitario muestran de forma consistente que la mayoría de estas personas combina una jornada laboral activa con las exigencias académicas nocturnas, lo que reduce el tiempo disponible para el trabajo autónomo y puede generar rezago progresivo \parencite{etitc_bienestar}. Esta situación no queda registrada en el sistema Adviser V.11 hasta que la nota de un parcial ya está reprobada, momento en que la intervención del equipo de Bienestar suele ser demasiado tardía. Construir un sistema que detecte señales de riesgo antes del cierre del primer corte no es solo una mejora tecnológica: es una respuesta de equidad hacia una población que combina trabajo, familia y estudio con recursos limitados. El \PDI 2025-2032 de la \ETITC estableció como meta reducir la deserción en 2,5 puntos porcentuales \parencite{etitc2025b}; sin una herramienta que actúe de manera preventiva, esa meta es difícil de alcanzar con la capacidad operativa actual del área de Bienestar. La presente investigación responde directamente a esa necesidad institucional desde un enfoque técnico.

Desde el punto de vista técnico, el proyecto resuelve dos problemas al mismo tiempo. El primero es la ausencia de datos individuales reales disponibles para el entrenamiento del modelo, que se abordó generando 1\,000 registros sintéticos calibrados con las distribuciones históricas de la institución, lo que permitió validar el sistema sin exponer información confidencial \parencite{ley1581}. El segundo es la desconexión entre los tres sistemas de información de la universidad —Adviser, Microsoft Teams y las encuestas de Bienestar—, resuelta mediante una arquitectura relacional de tres archivos \CSV unidos por un identificador único de estudiante. Los resultados obtenidos confirman que la arquitectura funciona: el modelo \RF clasificó correctamente el 95,0\,\% de los casos, y el modelo \XGB identificó el 85,7\,\% de los estudiantes en riesgo. El análisis de importancia de variables reveló, además, que la autoevaluación y la coevaluación del Corte~3 son los predictores más informativos del modelo, superando incluso a la nota definitiva del mismo corte, lo que justifica incorporarlas de forma sistemática en los registros docentes e instrumentos de Bienestar.

\newpage

% ──────────────────────────── HIPÓTESIS ──────────────────────────────────
\section{Hipótesis}

La integración de variables de caracterización socioeconómica inicial con variables de rendimiento académico estructuradas por los tres cortes institucionales, en una arquitectura relacional de datos sintéticos, permite identificar patrones de riesgo de deserción con mayor precisión que los métodos de seguimiento actualmente disponibles en la ETITC.

\newpage

% ──────────────────────────── METODOLOGÍA ────────────────────────────────
\section{Metodología}

\subsection{Tipo y enfoque de investigación}

Esta investigación es de tipo \textbf{aplicada}, con enfoque \textbf{cuantitativo y carácter predictivo}. Se parte del análisis descriptivo de los datos históricos de la institución para construir un modelo matemático que clasifique a los estudiantes según su nivel de riesgo de abandono. Al no contar con microdatos reales disponibles, la investigación incluye una fase de \textbf{simulación estadística con reglas de negocio} como método de validación controlada, estrategia coherente con el nivel de madurez tecnológica TRL~3 en que se ubica actualmente el proyecto.

\subsection{Vigilancia tecnológica}

La vigilancia tecnológica es el proceso sistemático de identificar, analizar y aprovechar las tendencias tecnológicas relevantes en un campo, con el fin de orientar mejor las decisiones de desarrollo \parencite{escorsa2003}. Este concepto fue desarrollado en el ámbito de la gestión de la innovación empresarial por Pere Escorsa y Jaume Valls, investigadores de la Universidad Politécnica de Cataluña, cuya obra es referencia estándar en el campo de la vigilancia tecnológica en habla hispana. En este proyecto se aplicó su metodología para revisar los sistemas de alerta temprana existentes a nivel nacional e internacional, con los siguientes resultados:

\begin{itemize}
  \item \textbf{SPADIES 3.0 (Colombia):} Sistema del Ministerio de Educación Nacional que calcula tasas de deserción a nivel institucional y nacional. Usa variables de matrícula y graduación, pero no integra datos de comportamiento académico continuo ni indicadores socioeconómicos individuales \parencite{men2025}.
  \item \textbf{Adviser V.11 y RUSIA (ETITC):} Plataformas internas que registran calificaciones y datos de matrícula.\parencite{bersoft2024}.
  \item \textbf{Herramientas internacionales de \EDM:} Investigaciones en minería de datos educativos demuestran que los modelos basados en registros de actividad continua en plataformas digitales —entregas de tareas, tiempos de respuesta, frecuencia de acceso— superan en capacidad predictiva a los modelos basados únicamente en calificaciones \parencite{macfadyen2010}. Sistemas comerciales como EAB Navigate y Civitas Learning en Estados Unidos aplican este enfoque, pero no están adaptados al contexto de la formación nocturna colombiana ni a la estructura de tres cortes institucionales.
  \item \textbf{Oportunidad de potenciación:} La revisión de los sistemas documentados sugiere que ninguno de ellos integra de forma explícita la caracterización socioeconómica con el comportamiento académico continuo por corte en un único modelo predictivo. Esta observación —que podría revisarse con una búsqueda más exhaustiva— orientó el diseño de la arquitectura propuesta como un complemento que busca potenciar las capacidades ya existentes de Adviser y RUSIA, no reemplazarlas.
\end{itemize}

\subsection{Modelo TRL}

La escala \TRL clasifica las tecnologías desde la idea básica (TRL~1) hasta el producto implementado y verificado en producción (TRL~9). La progresión del proyecto fue la siguiente:

\begin{itemize}
  \item \textbf{TRL~1 — Principios básicos observados:} Revisión de la literatura sobre \EDM y teoría de deserción; identificación de las variables relevantes en los informes de Bienestar \parencite{tinto1975, macfadyen2010}.
  \item \textbf{TRL~2 — Concepto tecnológico formulado:} Diseño de la arquitectura de tres archivos \CSV, definición de las reglas de negocio y selección de los algoritmos (\RF y \XGB).
  \item \textbf{TRL~3 — Demostración experimental (nivel actual):} Generación del conjunto de datos sintéticos, entrenamiento de los modelos y obtención de resultados cuantificables. El modelo \RF alcanzó exactitud del 95,0\,\% y AUC de 0,976; el modelo \XGB alcanzó sensibilidad del 85,7\,\%.
  \item \textbf{TRL~4 (próximo paso):} Integración con datos reales anonimizados de la institución en un entorno de prueba controlado, sujeto a acuerdo formal con la Dirección de Bienestar Universitario.
\end{itemize}

\subsection{Reglas de negocio del simulador}

Las reglas de negocio son las condiciones lógicas que determinaron cómo se generó cada registro sintético. Se calibraron con los promedios históricos de los informes de Bienestar \parencite{etitc_bienestar}.

En la notación empleada, \textbf{R1 a R4} identifican cada regla de simulación (R = Regla). \textbf{C1, C2 y C3} corresponden al \textbf{Corte~1, Corte~2 y Corte~3}, que son los tres períodos de evaluación parcial que estructuran el semestre académico en la ETITC: el primero cubre aproximadamente las primeras cinco semanas, el segundo las semanas seis a diez, y el tercero las semanas once al cierre del semestre.

\begin{enumerate}
  \item \textbf{Perfil base:} El 68\,\% de los registros tienen la condición de trabajador activo. El 45\,\% reporta debilidad en matemáticas. La tasa de deserción simulada resultó del 17,6\,\%, coherente con el rango histórico de la institución.

  \item \textbf{R1 — Presión laboral sobre asistencia:} Las inasistencias por corte se modelaron como una distribución normal truncada en el rango $[0,\, 8]$. Los estudiantes trabajadores reciben medias de $(3{,}0;\; 3{,}0;\; 4{,}0)$ para los cortes C1, C2 y C3, reflejando la fatiga acumulada hacia el final del semestre; los no trabajadores reciben medias de $(0{,}8;\; 0{,}8;\; 1{,}0)$.

  \item \textbf{R2 — Riesgo cognitivo en parciales:} Los estudiantes con más de tres años fuera del colegio y debilidad en matemáticas reciben una calificación media de 2,2 en parciales de ciencias básicas. El resto recibe media 3,5 (con presión laboral) o 4,2 (sin presión).

  \item \textbf{R3 — Barrera tecnológica en entregas digitales:} Los estudiantes que solo cuentan con un teléfono móvil como herramienta de trabajo reciben una penalización de $-0{,}5$ puntos en las entregas de Teams de cada corte, independientemente de su asistencia.

  \item \textbf{R4 — Abandono tácito:} Un estudiante se marca como desertor (\texttt{deserto~=~1}) si su promedio de notas definitivas baja de 3,0, o si registra autoevaluación y coevaluación en cero durante los cortes 2 o 3. La probabilidad de activar esta condición es del 50\,\% en C1, 70\,\% en C2 y 90\,\% en C3, representando la deserción progresiva tardía.
\end{enumerate}

\subsection{Procedimiento paso a paso}

\begin{enumerate}[label=\textbf{Fase \arabic*.}]
  \item \textbf{Revisión documental y definición de variables:} Se analizaron los informes de Bienestar Universitario de los últimos diez semestres para extraer las distribuciones de las variables con mayor peso en la deserción \parencite{etitc_bienestar}. Se definieron 24 variables distribuidas en tres archivos: caracterización socioeconómica, notas por corte (parciales, Teams, autoevaluación, coevaluación e inasistencias) y variable objetivo.

  \item \textbf{Generación de datos sintéticos:} Se programó un script en Python que aplicó las cuatro reglas de negocio para generar 1\,000 registros. El conjunto resultante tiene una tasa de deserción del 17,6\,\% (176 de 1\,000 estudiantes), coherente con los datos históricos.

  \item \textbf{Ingeniería de características y unión relacional:} Los tres archivos se unieron mediante el campo \texttt{id\_estudiante} en un proceso \ETL. Se aplicó partición estratificada del conjunto de datos: 80\,\% entrenamiento (800 registros) y 20\,\% prueba (200 registros).

  \item \textbf{Entrenamiento y evaluación de modelos:} Se entrenaron dos clasificadores: \RF y \XGB. Como el conjunto de datos tiene muchos más estudiantes sin deserción que con deserción (proporción aproximada de 5 a 1), se configuró el \XGB para que penalice más los casos no detectados que las alertas falsas; esto se logra con el parámetro \texttt{scale\_pos\_weight}~=~4,68, cuyo valor resulta de dividir el número de estudiantes sin deserción entre el número de desertores en el conjunto de entrenamiento \parencite{chen2016}. Los dos modelos se midieron con cinco indicadores: exactitud, precisión, sensibilidad, F1 y AUC-ROC. Al modelo con mejores resultados se le aplicó validación cruzada de cinco particiones (\textit{5-fold cross-validation}), que consiste en entrenar y evaluar el modelo cinco veces sobre subconjuntos distintos del conjunto de datos para confirmar que el resultado no depende de una única división favorable de los datos.
\end{enumerate}

\newpage

% ──────────────────────────── MARCOS ─────────────────────────────────────
\section{Marcos de investigación}

\subsection{Marco conceptual}

Los siguientes conceptos son fundamentales para entender el funcionamiento del sistema propuesto:

\begin{itemize}
  \item \textbf{Pobreza de tiempo:} Escasez severa de horas disponibles para el autoestudio, causada por la superposición de una jornada laboral activa con las exigencias académicas nocturnas \parencite{arias2018}. En la ETITC, más del 52\,\% de los estudiantes nocturnos declara este problema de forma recurrente \parencite{etitc_bienestar}.

  \item \textbf{Huella digital académica:} El rastro medible que genera un estudiante al interactuar con plataformas institucionales. En este proyecto se mide como la frecuencia y puntualidad de las entregas en Microsoft Teams por corte.

  \item \textbf{Deserción temprana:} Abandono antes del cierre del primer corte académico.

  \item \textbf{Deserción tardía:} Retiro silencioso que ocurre entre el segundo y el tercer corte, generalmente por acumulación de rezago y fatiga. Es la forma más frecuente entre los estudiantes nocturnos de la ETITC \parencite{etitc_bienestar}.

  \item \textbf{Datos sintéticos relacionales:} Registros generados por un algoritmo que replica las distribuciones y correlaciones de una población real sin exponer información confidencial de individuos reales.

  \item \textbf{Variable objetivo:} En \ML supervisado, la variable que el modelo aprende a predecir. En este proyecto es \texttt{deserto} (1 = desertó, 0 = continuó).

  \item \textbf{Autoevaluación y coevaluación:} Notas de la componente cualitativa del sistema de evaluación institucional. En el análisis de importancia de variables, estas dos señales sumadas representaron el 20,7\,\% del poder predictivo del modelo, superando a la nota definitiva del Corte~3 (18,9\,\%).

  \item \textbf{Falso negativo (FN):} Estudiante en riesgo que el modelo no identificó. En el contexto de la deserción estudiantil, un FN significa un caso que queda sin atención de Bienestar, lo cual tiene consecuencias directas sobre la retención.

  \item \textbf{Falso positivo (FP):} Alerta emitida sobre un estudiante que en realidad no estaba en riesgo. Genera trabajo innecesario para el equipo de Bienestar.
\end{itemize}

\subsection{Marco teórico}

La investigación se fundamenta en tres corrientes teóricas que explican la deserción desde ángulos complementarios:

\begin{enumerate}
  \item \textbf{Teoría de la integración académica y social:} \textcite{tinto1975} demostró que la permanencia de un estudiante en la educación superior depende de su grado de integración con las estructuras académicas y sociales de la institución. Cuando esa integración se debilita —por carga laboral, fatiga acumulada o rezago curricular— la deserción es una consecuencia previsible. Este modelo explica por qué las variables de autoevaluación y coevaluación resultaron ser los predictores más potentes: miden directamente el nivel de vinculación activa del estudiante con su entorno académico.

  \item \textbf{Minería de datos educativos (\EDM):} Tal como lo plantean \textcite{romero2010}, la \EDM consiste en aplicar técnicas de análisis de datos para extraer patrones útiles de los entornos de aprendizaje. Dentro de este campo, \textcite{macfadyen2010} demostraron que los registros de actividad en plataformas \LMS —especialmente la entrega oportuna de tareas— tienen mayor poder predictivo que el simple acceso pasivo a materiales. Este hallazgo respalda la inclusión de las notas de Teams y los conteos de inasistencia como variables centrales del modelo propuesto.

  \item \textbf{Algoritmos supervisados para clasificación:} Los dos modelos empleados tienen fundamentos teóricos bien establecidos. El \RF, propuesto por \textcite{breiman2001}, construye un conjunto de árboles de decisión independientes cuyas predicciones se combinan para reducir el sobreajuste y mejorar la estabilidad del resultado. El \XGB, desarrollado por \textcite{chen2016}, construye los árboles de forma secuencial: cada nuevo árbol corrige los errores del anterior. La ventaja del \XGB en este proyecto radica en el parámetro \texttt{scale\_pos\_weight}, que indica al modelo cuánto más grave es dejar pasar un desertor sin detectar que emitir una alerta incorrecta.
\end{enumerate}

\subsection{Marco normativo}

El diseño del sistema se rige por la legislación colombiana vigente en materia de protección de datos personales y uso de tecnologías:

\begin{itemize}
  \item \textbf{Ley Estatutaria 1581 de 2012 \parencite{ley1581}:} Establece las reglas para el tratamiento de datos personales en Colombia. Por tratarse de información altamente sensible —calificaciones y condiciones socioeconómicas de estudiantes—, el prototipo opera con identificadores cifrados anonimizados (\texttt{id\_estudiante}) en lugar de nombres o documentos de identidad.

  \item \textbf{Decreto 1377 de 2013 \parencite{decreto1377}:} Reglamenta parcialmente la Ley 1581. Define los esquemas de autorización y las políticas de tratamiento de datos que la institución deberá implementar cuando decida alimentar el modelo con información real de sus estudiantes.

  \item \textbf{Ley 1341 de 2009 (Ley TIC) \parencite{ley1341}:} Promueve el acceso y uso de las tecnologías de la información para el desarrollo social. Respalda la pertinencia de construir herramientas de inteligencia artificial orientadas a la permanencia estudiantil en el contexto colombiano.
\end{itemize}


% ──────────────────────────── TIPOS DE DATOS ────────────────────────────
\section{Tipos de datos}

El sistema procesó cuatro categorías de datos distribuidas en los tres archivos \CSV relacionales:

\begin{itemize}
  \item \textbf{Datos categóricos nominales y ordinales:} Procedentes del archivo de caracterización: condición laboral (\texttt{trabaja}: sí/no), tiempo desde el bachillerato (\texttt{bachi\_tiempo}: menos de 1 año, de 1 a 3 años, más de 3 años), estrato socioeconómico (1 a 6) y tipo de dispositivo tecnológico disponible \parencite{etitc_bienestar}.

  \item \textbf{Datos numéricos continuos:} Calificaciones de parciales presenciales, notas de entregas en Teams y notas de autoevaluación y coevaluación, todas en escala de 0,0 a 5,0.

  \item \textbf{Datos numéricos discretos:} Conteo de sesiones no asistidas por corte (\texttt{inasistencias\_cX}), con rango entero $[0,\, 8]$, que permite capturar la presión laboral de manera proporcional en lugar de un simple indicador de presencia o ausencia.

  \item \textbf{Datos booleanos:} Indicadores binarios para barrera tecnológica (\texttt{barrera\_tecnologica}), debilidad en matemáticas (\texttt{debilidad\_mate}) y la variable objetivo (\texttt{deserto}: 1 = abandonó, 0 = continuó).
\end{itemize}

% ─────────────────── RECOLECCIÓN DE DATOS CUANTITATIVOS ─────────────────
\section{Recolección de datos cuantitativos e instrumentos}

\subsection*{Instrumento principal: encuesta de caracterización}

El instrumento de recolección cuantitativa es el formulario de caracterización estudiantil que aplica Bienestar Universitario al inicio de cada semestre \parencite{etitc_bienestar}. Este formulario recoge variables socioeconómicas como: condición laboral, tiempo de desvinculación académica desde el bachillerato, estrato, situación de vivienda, acceso a dispositivos tecnológicos y percepción de dificultades en ciencias básicas. Esta información constituye el archivo \texttt{bienestar\_caracterizacion.csv} de la arquitectura propuesta.

\subsection*{Fuente complementaria: planillas de notas}

Las calificaciones por corte se obtienen del sistema Adviser V.11 \parencite{bersoft2024}, que registra las notas de parciales presenciales al cierre de cada uno de los tres cortes institucionales. Las notas de entregas de actividades, de autoevaluación y de coevaluación se obtienen de los canales de la asignatura en Microsoft Teams. Estos datos conforman el archivo \texttt{adviser\_teams\_notas.csv}.

\subsection*{Datos generados en esta fase}

Al no contar con acceso a los microdatos reales, se generó un conjunto sintético de 1\,000 registros con 24 variables distribuidas en los tres archivos mencionados. La variable objetivo (\texttt{deserto}) se almacena en el archivo \texttt{registro\_admisiones\_target.csv}. La tasa de deserción resultante fue del 17,6\,\% (176 registros positivos de un total de 1\,000).

% ──────────────────── RECOLECCIÓN DE DATOS CUALITATIVOS ─────────────────
\section{Recolección de datos cualitativos}

El componente cualitativo se captura a través de dos variables formalizadas en el sistema de evaluación institucional:

\begin{itemize}
  \item \textbf{Autoevaluación:} Nota que el propio estudiante asigna a su desempeño al cierre de cada corte.
  \item \textbf{Coevaluación:} Nota asignada por los compañeros de equipo al trabajo colaborativo del estudiante.
\end{itemize}

Aunque son de naturaleza cualitativa —reflejan percepciones y compromisos—, ambas se expresan en la planilla docente como calificaciones numéricas en escala de 0,0 a 5,0, lo que permite incluirlas directamente en el modelo cuantitativo. Los resultados del análisis de importancia de variables confirmaron que la autoevaluación del Corte~3 es el segundo predictor más importante (13,1\,\%) y la coevaluación del mismo corte el sexto (7,6\,\%). Sumadas, representan el 20,7\,\% del poder predictivo del modelo, lo que valida su uso como sensores tempranos de desconexión académica.

\newpage

% ─────────────────── ANÁLISIS DE RESULTADOS ──────────────────────────────
\section{Análisis de resultados cuanti-cualitativos}

\subsection{Resultados del modelo Random Forest}

El modelo \RF entrenado sobre el conjunto de datos sintéticos, con umbral de clasificación 0,5, obtuvo los siguientes resultados sobre el conjunto de prueba de 200 registros (35 desertores, 165 no desertores):

\begin{itemize}
  \item \textbf{Exactitud (Accuracy):} 95,0\,\% — el modelo clasificó correctamente 190 de 200 estudiantes.
  \item \textbf{Precisión (Precision):} 100\,\% — ninguna de las alertas emitidas fue incorrecta.
  \item \textbf{Sensibilidad (Recall):} 71,4\,\% — el modelo identificó 25 de los 35 estudiantes en riesgo.
  \item \textbf{Puntaje F1:} 83,3\,\% — balance entre precisión y sensibilidad.
  \item \textbf{Área bajo la curva ROC (AUC):} 0,976 — capacidad discriminativa muy alta. El AUC mide qué tan bien el modelo separa los estudiantes en riesgo de los que no lo están; su rango va de 0,5 (equivalente a adivinar al azar) a 1,0 (discriminación perfecta). Un valor de 0,976 significa que si se toman al azar un desertor y un no desertor del conjunto de prueba, el modelo asigna mayor probabilidad de deserción al desertor correcto en el 97,6\,\% de los casos.
  \item \textbf{Falsos negativos (FN):} 10 — estudiantes en riesgo que el modelo no detectó.
  \item \textbf{Falsos positivos (FP):} 0 — ninguna alerta incorrecta.
\end{itemize}

Al reducir el umbral de clasificación de 0,5 a 0,35, el modelo detecta 4 FN menos (pasa de 10 a 6 casos no detectados) a costa de 3 alertas incorrectas.

\subsection{Resultados del modelo XGBoost}

El modelo \XGB, entrenado con \texttt{scale\_pos\_weight}~=~4,68 para compensar el desequilibrio entre las clases, obtuvo los siguientes resultados con umbral 0,5 \parencite{chen2016}:

\begin{itemize}
  \item \textbf{Exactitud:} 95,5\,\%.
  \item \textbf{Precisión:} 88,2\,\% — 4 de las 34 alertas emitidas fueron incorrectas.
  \item \textbf{Sensibilidad:} 85,7\,\% — identificó 30 de los 35 estudiantes en riesgo.
  \item \textbf{Puntaje F1:} 87,0\,\%.
  \item \textbf{AUC-ROC:} 0,973.
  \item \textbf{FN:} 5 — reduce a la mitad los casos no detectados respecto al RF.
  \item \textbf{FP:} 4 — cuatro alertas incorrectas por cada 200 estudiantes.
\end{itemize}

\subsection{Importancia de variables}

Las diez variables de mayor peso predictivo en el modelo \RF, ordenadas de mayor a menor, fueron:

\begin{enumerate}
  \item Nota definitiva Corte~3: 18,9\,\%
  \item Autoevaluación Corte~3: 13,1\,\%
  \item Nota definitiva Corte~2: 10,6\,\%
  \item Inasistencias Corte~3: 8,9\,\%
  \item Nota definitiva Corte~1: 7,8\,\%
  \item Coevaluación Corte~3: 7,6\,\%
  \item Parcial Corte~2: 5,3\,\%
  \item Parcial Corte~1: 4,6\,\%
  \item Notas Teams Corte~3: 4,5\,\%
  \item Parcial Corte~3: 3,7\,\%
\end{enumerate}

El hallazgo más relevante es que la autoevaluación y la coevaluación del Corte~3 sumadas aportan el 20,7\,\% del poder predictivo del modelo, superando a la nota definitiva del mismo corte.

\newpage

% ──────────────────── PROPUESTA DE PROTOTIPO Y TRL ───────────────────────
\section{Propuesta de prototipo basado en análisis y modelo TRL}

El prototipo desarrollado se ubica actualmente en \textbf{TRL~3}, nivel que corresponde a la demostración experimental de la función la tecnología. La arquitectura técnica consiste en tres archivos \CSV relacionales unificados mediante un proceso \ETL programado en Python:

\begin{itemize}
  \item \texttt{bienestar\_caracterizacion.csv}: Perfil socioeconómico inicial (condición laboral, estrato, tiempo desde el bachillerato, barrera tecnológica, debilidad en matemáticas).
  \item \texttt{adviser\_teams\_notas.csv}: Rendimiento académico continuo — parciales, entregas de Teams, autoevaluación, coevaluación e inasistencias para cada uno de los tres cortes.
  \item \texttt{registro\_admisiones\_target.csv}: Variable objetivo (\texttt{deserto}), derivada aplicando las reglas de negocio sobre los archivos anteriores.
\end{itemize}

El flujo de procesamiento (\textit{pipeline}) de \ML fue implementado en Python con las librerías \textit{Scikit-learn} y \textit{XGBoost}. Un \textit{pipeline} es una secuencia automatizada de pasos de transformación y modelado encadenados, de modo que los datos pasan de un paso al siguiente sin intervención manual entre etapas, reduciendo errores y facilitando la reproducibilidad. La arquitectura del script se organiza en cuatro módulos:

\begin{enumerate}
  \item \textbf{Carga y unión relacional:} Los tres archivos \CSV se cargan con la librería \textit{pandas} y se integran mediante el campo \texttt{id\_estudiante}, que actúa como llave foránea entre los tres conjuntos de datos, siguiendo el modelo relacional descrito en la sección anterior.
  \item \textbf{Preprocesamiento:} Las variables numéricas se escalan al rango $[0,\,1]$ (\texttt{MinMaxScaler}) para que ninguna variable domine al modelo por su magnitud; las variables categóricas (como condición laboral o tiempo de desvinculación desde el bachillerato) se convierten a valores numéricos (\texttt{OrdinalEncoder}).
  \item \textbf{Entrenamiento y evaluación:} Se entrena primero el \RF y luego el \XGB con los parámetros descritos. Cada modelo se evalúa sobre el 20\,\% de datos reservados para prueba (200 registros).
  \item \textbf{Análisis de importancia de variables:} Se extrae el atributo \texttt{feature\_importances\_} del modelo ganador, que indica qué porcentaje del poder predictivo aportó cada variable al resultado final.
\end{enumerate}

Los resultados obtenidos confirman la viabilidad técnica de la arquitectura.

\textbf{Hoja de ruta hacia TRL~4:}
\begin{enumerate}
  \item Ajustar el umbral del modelo \RF a 0,35 para reducir los casos no detectados de 10 a 6, con un costo de 3 alertas incorrectas adicionales.
  \item Obtener el aval formal de la Dirección de Bienestar Universitario y la Facultad de Sistemas para una prueba piloto con datos reales anonimizados.
  \item Validar el modelo entrenado sobre datos sintéticos con información real de al menos dos semestres, para verificar que la precisión se mantiene fuera del entorno simulado.
  \item Diseñar una interfaz de alertas para el personal de Bienestar que explique, por cada estudiante en riesgo, cuáles son las variables críticas que activaron la alerta.
\end{enumerate}

\newpage

% ──────────────────────────── CONCLUSIONES ───────────────────────────────
\section{Conclusiones}

\begin{enumerate}
  \item La arquitectura relacional de tres fuentes (\CSV de caracterización, notas por corte y variable objetivo) demostró ser técnicamente viable para replicar la estructura fragmentada de los sistemas de información de la ETITC y entrenar clasificadores de \ML con datos calibrados a las distribuciones históricas de la institución \parencite{etitc_bienestar}.

  \item El modelo \RF alcanzó una exactitud del 95,0\,\% y un AUC de 0,976 sobre los datos sintéticos, sin generar ninguna alerta incorrecta. El modelo \XGB mejoró la detección de estudiantes en riesgo (sensibilidad del 85,7\,\% frente al 71,4\,\% del RF), a costa de cuatro alertas incorrectas por cada 200 estudiantes evaluados. Esta diferencia es operacionalmente relevante: para una cohorte de 500 estudiantes nocturnos, el XGB dejaría sin identificar aproximadamente la mitad de los casos que el RF no detecta.

  \item La autoevaluación y la coevaluación del Corte~3 sumadas representaron el 20,7\,\% del peso predictivo total del modelo, superando a la nota definitiva del mismo corte (18,9\,\%). Este hallazgo tiene una implicación directa para la práctica institucional: incorporar estas variables de forma sistemática en los registros docentes puede incrementar la capacidad de detección temprana incluso antes de que las notas finales queden consolidadas.

  \item El enfoque de simulación con reglas de negocio calibradas a datos históricos demostró ser un método válido para desarrollar y probar sistemas de \ML en entornos educativos donde los microdatos individuales no están disponibles por restricciones normativas \parencite{ley1581}. Esto abre una ruta metodológica replicable para otras instituciones colombianas en la misma situación.
\end{enumerate}

% ──────────────────────────── DISCUSIÓN ──────────────────────────────────
\section{Discusión}

Los resultados obtenidos son coherentes con la teoría de integración académica de \textcite{tinto1975}: las variables que mejor predicen la deserción no son las notas finales, sino los indicadores de desconexión progresiva —autoevaluación baja, inasistencias acumuladas, abandono de las entregas de Teams— que preceden al fracaso académico oficial semanas antes del cierre del corte. Esto confirma que el sistema propuesto actúa en el momento adecuado: antes, no después.

Sin embargo, es necesario reconocer las limitaciones del estudio. Los datos sobre los que se validó el modelo son sintéticos y fueron generados por el mismo equipo de investigación. Aunque las reglas de negocio se calibraron con los informes históricos de Bienestar \parencite{etitc_bienestar}, no es posible garantizar que todos los patrones del comportamiento humano real estén representados en la simulación. La deserción real responde también a eventos imprevistos —pérdida de empleo, enfermedad, cambio de ciudad— que las reglas lógicas no modelan de forma completa. Esta limitación es inherente a la fase TRL~3 y se resuelve en la fase TRL~4 al incorporar datos reales.

La comparación entre \RF y \XGB arroja una conclusión operativa clara. El \RF es preferible cuando la institución prioriza la ausencia de falsas alarmas, para no sobrecargar al equipo de Bienestar con casos innecesarios. El \XGB es preferible cuando la prioridad es no dejar pasar ningún estudiante en riesgo.

Frente al estado actual de las herramientas institucionales —Adviser V.11 y RUSIA—, que no generan ningún tipo de alerta predictiva automatizada \parencite{bersoft2024}, la brecha de valor que aportaría este sistema es significativa. El paso siguiente es obtener el respaldo institucional necesario para realizar una prueba piloto con datos reales, lo que llevaría el proyecto al nivel TRL~4 y permitiría cuantificar el impacto real sobre la retención estudiantil en la jornada nocturna.

\newpage
\printbibliography[title={Referencias bibliográficas}]

\end{document}
